%----------------------------------------------------------------------------------------
%	PACKAGES AND OTHER DOCUMENT CONFIGURATIONS
%----------------------------------------------------------------------------------------

\documentclass[a4paper,10pt]{article} % Default font size and paper size

\usepackage{fontspec} % For loading fonts
\defaultfontfeatures{Mapping=tex-text}
\setmainfont[SmallCapsFont = Fontin SmallCaps]{Fontin} % Main document font

\usepackage{xunicode,xltxtra,url,parskip} % Formatting packages

\usepackage{scrpage2} % Provides headers and footers configuration

\usepackage[usenames,dvipsnames]{xcolor} % Required for specifying custom colours

\usepackage{hyperref} % Required for adding links	and customizing them
\definecolor{linkcolour}{rgb}{0,0.2,0.6} % Link colour
\hypersetup{colorlinks,breaklinks,urlcolor=linkcolour,linkcolor=linkcolour} % Set link colours throughout the document

\usepackage{tabularx,colortbl} % Advanced table configurations
\usepackage{array} % Custom arrangement for the columns
\usepackage{multirow} %rows spanning multiple rows

\usepackage{titlesec} % Used to customize the \section command
\titleformat{\section}{\Large\scshape\raggedright}{}{0em}{}[\titlerule] % Text formatting of sections

\titlespacing{\section}{0pt}{3pt}{3pt} % Spacing around sections

\usepackage{enumitem} % For left margin removal in itemize

\usepackage{marvosym} % Allows the use of symbols

%----------------------------------------------------------------------------------------
% PAGE GEOMETRY
%----------------------------------------------------------------------------------------

\usepackage[top=1cm, bottom=0.5cm, lmargin=0.5cm, rmargin=1.3cm]{geometry} % The page Geometry

%----------------------------------------------------------------------------------------
% CUSTOM ENVIRONMENTS
%----------------------------------------------------------------------------------------
\newcommand{\work}[3]{
    \begin{tabular}{>{\raggedleft}p{2.2cm}|p{0.9\linewidth}}
        \textsc{#1} & \textcolor{NavyBlue}{#2}
                    \footnotesize{#3}
    \end{tabular}
}

\newcommand{\project}[4]{
    \begin{tabular}{p{0.85\linewidth}r}
        \textcolor{NavyBlue}{#2} & \multicolumn{1}{m{3cm}}{\raggedleft \textsc{#1}}\\
        #3
    \end{tabular}
    \begin{tabular}{p{\linewidth}}
        \footnotesize{#4}
    \end{tabular}
    \vspace{-0.1cm}
}

\newcommand{\lproject}[4]{
    \begin{tabular}{p{0.80\linewidth}r}
        \textcolor{NavyBlue}{#2} & \multicolumn{1}{m{4cm}}{\raggedleft \textsc{#1}}\\
        #3
    \end{tabular}
    \begin{tabular}{p{\linewidth}}
    \vspace{-0.3cm}
        \footnotesize{#4}
    \end{tabular}
    \vspace{-0.5cm}
}

\newcommand{\cproject}[5]{
    \begin{tabular}{p{0.80\linewidth}r}
        \textcolor{NavyBlue}{#2} & \multicolumn{1}{m{4cm}}{\raggedleft \textsc{#1}}\\
        #3 & #4
    \end{tabular}
    \begin{tabular}{p{\linewidth}}
    \vspace{-0.3cm}
        \footnotesize{#5}
    \end{tabular}
    \vspace{-0.5cm}
}

\newcommand{\scproject}[5]{
    \begin{tabular}{p{0.80\linewidth}r}
        \textcolor{NavyBlue}{#2} & \multicolumn{1}{m{4cm}}{\raggedleft \textsc{#1}}\\
        #3 & #4
    \end{tabular}
    \begin{tabular}{p{\linewidth}}
    \vspace{-0.3cm}
        \footnotesize{#5}
    \end{tabular}
    \vspace{-0.3cm}
}

\newcommand{\iproject}[3]{
    \begin{tabular}{p{0.85\linewidth}r}
        \textcolor{NavyBlue}{#2} & \multicolumn{1}{m{3cm}}{\raggedleft \textsc{#1}}\\
    \end{tabular}
    \begin{tabular}{p{\linewidth}}
    \vspace{-0.3cm}
        \footnotesize{#3}
    \end{tabular}
    \vspace{-0.5cm}
}

\newcommand{\siproject}[3]{
    \begin{tabular}{p{0.85\linewidth}r}
        \textcolor{NavyBlue}{#2} & \multicolumn{1}{m{3cm}}{\raggedleft \textsc{#1}}\\
    \end{tabular}
    \begin{tabular}{p{\linewidth}}
    \vspace{-0.3cm}
        \footnotesize{#3}
    \end{tabular}
    \vspace{-0.3cm}
}

\newcommand{\projectlist}[2]{
    \begin{tabular}{p{\linewidth}}
        \textcolor{NavyBlue}{#1}\\
        \vspace{-0.3cm}
        \footnotesize{#2}
    \end{tabular}
    \vspace{-0.4cm}
}

\newcommand{\itemlist}[1]{
    \begin{tabular}{>{\raggedleft}p{1.6cm}p{0.9\linewidth}}
        #1
    \end{tabular}
}

\newcommand{\skill}[2]{
    \begin{tabular}{p{0.85\linewidth}r}
        #2 & \multicolumn{1}{m{3cm}}{\raggedleft \textsc{#1}}\\
    \end{tabular}
    \vspace{-0.5cm}
}

\newcommand{\github}{
    \includegraphics[height=9pt]{icons/octa.png}
}

%----------------------------------------------------------------------------------------
% DOCUMENT BEGINS HERE
%----------------------------------------------------------------------------------------

\begin{document}

\font\fb=''[cmr10]'' % Change the font of the \LaTeX command under the skills section

\titleformat{\section}{\large\scshape\raggedright}{}{0em}{}[\titlerule] % Section formatting

%----------------------------------------------------------------------------------------
% HEADER SECTION
%----------------------------------------------------------------------------------------

{
    \begin{tabular}{rl}
        \multirow{3}{0.73\linewidth}{
            {
                \addfontfeature{LetterSpace=20.0}\fontsize{24}{24}\selectfont\scshape
                Srijan R. Shetty
            }\\
            \hspace{0.15cm} \emph{Department of Computer Science and Engineering}
        }

        & {\Large\Mobilefone} +(91) 9005900383 \\
        & {\Large\Letter} \href{mailto:srijan.shetty@gmail.com}{srijan.shetty@gmail.com}\\
        & {\github} \href{https://github.com/srijanshetty}{https://github.com/srijanshetty}\\
    \end{tabular}
}

%----------------------------------------------------------------------------------------
%   INTERESTS
%----------------------------------------------------------------------------------------
\section{Interests}

Programming Languages, Networks and Systems, Security, Web-Development, Natural Language Processing

%----------------------------------------------------------------------------------------
%	EDUCATION
%----------------------------------------------------------------------------------------
\section{Education}

\begin{tabular}{>{\raggedleft}p{1.5cm}p{15.5cm}r}

    \textsc{Current} & B. Tech in \textsc{Computer Science and Engineering},
                       \textbf{Indian Institute of Technology}, Kanpur
                     &   9.5/10.0\\

    \textsc{July 2011} & 12$^{th}$ Board, \textsc{CBSE} Board,
                         \normalsize\textbf{The Emerald Heights International School}, Indore
                       &   91.4\% \\

    \textsc{July 2011} & 10$^{th}$ Board, \textsc{ICSE} Board,
                         \normalsize\textbf{The Laurels School International}, Indore
                       &   95.4\% \\

\end{tabular}

%----------------------------------------------------------------------------------------
%	SCHOLARSHIPS AND CONFERENCES
%----------------------------------------------------------------------------------------
\section{Scholarships and Conferences}


\itemlist {
    \textsc{Jan 2011} & Kishore Vaigyanik Protsahan Yojana (KVPY) Fellowship \\
    \textsc{Dec 2009} & CSIR Programme on Youth for Leadership in Science held at Advanced Materials Research Institute, Bhopal\\
}


%----------------------------------------------------------------------------------------
%   SCHOLASTIC ACHIEVEMENTS
%----------------------------------------------------------------------------------------
\section{Scholastic Achievements}

\itemlist {
    \textsc{2011-2014}   & Received an \textbf{A* grade}, for exceptional performance in
                            \textbf{Computer, Internet and Network Security};
                            \textbf{Computer Networks};
                            \textbf{Computer Organization};
                            \textbf{Logic for Computer Science};
                            \textbf{Fundamentals of Computing};
                            \textbf{Introduction to Philosophy} \\
    \textsc{2012-2013}   & \textbf{Academic Excellence Award}, IIT Kanpur. \\
    \textsc{2011-2012}   & \textbf{Academic Excellence Award}, IIT Kanpur. \\
    \textsc{May 2011}    & Secured All India rank \textbf{1937} in Joint Entrance Exam (JEE) \\
    \textsc{Jan 2010}    & Qualified \textbf{Regional Mathematics Olympiad} (RMO) \\
    \textsc{2009, 2005}  & Awarded \textbf{Overall Best Student} for displaying all round excellence. \\
    \textsc{Dec 2008}    & Secured All India Rank 297 in FIITJEE Talent Reward Examination. \\
    \textsc{2000-2009}   & Recipient of \textbf{Scholastic Excellence} awards for exemplary academic performance throughout
                            2$^{nd}$ to 12$^{th}$ grade. \\
}

%----------------------------------------------------------------------------------------
%   SKILL SET
%----------------------------------------------------------------------------------------
\section{Skill Set}

\itemlist {
    \textsc{Languages} %
            & \textbf{JavaScript} (Expert), \textbf{Python} (Proficient), \textbf{C++} (Proficient), \textbf{C\#} (Efficient),
              \textbf{C} (Proficient)\\
    \textsc{Web Dev} %
            & \textbf{HTML5}, \textbf{CSS}, \textbf{SQL and NoSQL}, \textbf{AngularJS}, \textbf{Nodejs}\\
    \textsc{Tools} %
            & \textbf{Git} (Expert), \textbf{Shell Scripting} (Proficient), \textbf{Prezi} (Proficient), \textbf{Vim} (Expert), \textbf{LaTeX}
              (Proficient)\\
}

%----------------------------------------------------------------------------------------
%   WORK EXPERIENCE
%----------------------------------------------------------------------------------------

\section{Work Experience}

\work {May-July 2014}
      {Research Intern at \textsc{Microsoft Research India}, Bangalore}
      {
         \begin{itemize}[leftmargin=*]
             \item Analysed the implementations of different distributed systems and databases architectures like
                 \textbf{Google MapReduce}, \textbf{Google FileSystem}, \textbf{Google MillWheel} \textbf{Yahoo! PNUTS},
                 \textbf{Amazon DynamoDB}, \textbf{Apache Kafka}, \textbf{Apache Spark}, \textbf{Apache Hadoop}, \textbf{Apache Storm},
                 \textbf{Microsoft Dryad}, and \textbf{Microsoft DryadLINQ} to understand the design decisions involved.
             \item Used \textbf{TPL Dataflow} and \textbf{Lightweight Tasks} to implement concurrent pipeline processing of messages in
                 CScale to preserve processing order of partially ordered messages.
             \item Performed \textbf{performance testing} of the concurrency optimizations in real world scenarios leading to a gain in processing time.
             \item Performed \textbf{failure testing} of the existing infrastructure and debugged detrimental replication errors in the
                 existing infrastructure.
         \end{itemize}
     }

\work {May-July 2013}
      {Software Development at \textsc{Aurus Networks'}, Bangalore}
      {
          \begin{itemize}[leftmargin=*]
              \item Ported the legacy flash based web player in \textbf{CourseHub} --- Aurus Networks'
                  MOOC offering --- to a technology agnostic \textbf{HTML5/JavaScript} web player leading
                  to a uniform cross-platform viewing experience.
              \item Developed a \textbf{XMPP} based \textbf{real-time chat service} using Jabber framework,
                  for conversations during live lectures.
              \item Engineered an extensible Jabber framework for real-time communications between different services.
          \end{itemize}
      }

%----------------------------------------------------------------------------------------
%  COURSE PROJECTS
%----------------------------------------------------------------------------------------
\let\thefootnote\relax\footnote{Monsoon semester: July to November}
\let\thefootnote\relax\footnote{Winter Semester: January to April}
\section{Course Projects}

\cproject
    {Monsoon 2014-2015}
    {Multi Factor Authentication in OpenVPN}
    {\textsc{\raggedright Mozilla Winter of Security}, Guillaume Destuynder and Professor Dheeraj Sanghi}
    {\href{https://wiki.mozilla.org/Security/Mentorships/MWoS/2014/OpenVPN\_MFA}{ \footnotesize{[Project Wiki]}} }
    {
      \begin{itemize}[leftmargin=0.5cm]
          \item One of the eleven projects to be selected as a part of Mozilla's Winter of Security Initiative 2014.
          \item The objective of the project is to implement true arbitrary multi-factor authentication support in
              OpenVPN; and to implement session-support for the additional factors.
      \end{itemize}
    }

\scproject
    {Winter 2013-2014}
    {JavaScript to MIPS Compiler}
    {\textsc{\raggedright Compilers}, Professor Subhajit Roy}
    { \href{https://github.com/srijanshetty/javascript-compiler} { \footnotesize{[Repository]}} }
    {
          Implemented an end-to-end ECMAScript 5.1 compiler for the MIPS architecture with support for first class
          functions, anonymous functions, static types, type annotations, recursion and nested functions.
    }

\cproject
    {Winter 2013-2014}
    {Hindi author attribution}
    {\textsc{\raggedright Artificial Intelligence}, Professor Amitabha Mukherjee}
    { \href{https://github.com/srijanshetty/author-attribution}{ \footnotesize{[Repository]}} }
    {
     \begin{itemize}[leftmargin=0.5cm]
         \item Presented a poster on Opinion Word Expansion and Target Extraction through Double Propogation, Guang Qin,
             Bing Liu, Jiajun Bu, Chun Chen that was nominated as the best poster by all students.
         \item Selected as one of the top seven projects at the end of the term.
         \item Generated a Hindi literature corpus consisting of works by Rabindranath Tagore,
             Vibhuti Narayan, Premchand, Sarat Chandra Chattopadhyay and Dhamarvir Bharati.
         \item Clustered a sample set of multiple Hindi authors using K-means unsupervised clustering to achieve
             F-scores in the range of 90-97\% for each author.
         \item Classified each author using Support Vector Machine Classification using Radial Basis Kernel function
             and achieved an F-score of 90--95\% for each author.
         \item Performed Multiple Discriminant Analysis on the data for a quantitative comparison with word
             collocations; there was no appreciable difference in the results with the F-scores increasing only
             marginally.
     \end{itemize}
    }

\cproject
    {Monsoon 2013-2014}
    {NachOS}
    {\textsc{Operating Systems}, Professor Mainak Chaudhari}
    {}
    {
     \begin{itemize}[leftmargin=0.5cm]
         \item Extended the standard system call library of NachOS and implemented system calls pertaining to Fork, Exec,
             Join, Yield, Sleep and Exit.
             \href{https://github.com/srijanshetty/nachos-syscalls/}{[Repository]}
         \item Implemented process scheduling algorithms: UNIX Scheduling, First in First Out,
             Round Robin, Shortest Job First and Non-pre-emptive job scheduling to assess their relative performances.
             \href{https://github.com/srijanshetty/nachos-scheduling}{[Repository]}
         \item Programmed page replacement algorithms: Random Page Allocation, First in First Out,
             Least Recently Used(LRU) and LRU Clock to evaluate relative performances under difference scenarios.
             \href{https://github.com/srijanshetty/nachos-final/}{[Repository]}
      \end{itemize}
    }

\cproject
    {Monsoon 2013-2014}
    {P2P File Sharing and Streaming and IP Spoofing}
    {\textsc{Computer Networks}, Professor Dheeraj Sanghi}
    { \href{https://github.com/srijanshetty/nodesock}{ \footnotesize{[Repository]} } }
    {
        \begin{itemize}[leftmargin=0.5cm]
            \item Conceived and implemented a peer-to-peer protocol for file and media transfer in Node.js to share
                data over any LAN.
            \item Generated raw IP packets with spoofed IP addresses to exploit test machines with ARP Poisoning,
                Smurfing and SYN Flooding.
        \end{itemize}
    }

\cproject
    {Monsoon 2013-2014}
    {Prolog Shortest Paths and Erlang Traffic Server}
    {\textsc{Principles of Programming Languages}, Professor Piyush Kurur}
    {}
    {
       \begin{itemize}[leftmargin=0.5cm]
           \item Implemented shortest path algorithm in Prolog using Logic Programming on the Delhi Metro route.
           \item Implemented a simple traffic server in Erlang using message passing concurrency and tail call recursion.
       \end{itemize}
    }

\cproject
    {Winter 2012-2013}
    {8-bit General Purpose Computer on a FPGA}
    {\textsc{Computer Organization}, Professor Subhajit Roy}
    { \href{https://github.com/srijanshetty/220\_y11} {\footnotesize{[Repository]}} }
    {
      \begin{itemize}[leftmargin=0.5cm]
          \item Designed and programmed a simple Instruction set Architecture (ISA) for a 8-bit General Purpose Computer with a load-store
              architecture on \textbf{Xilinx Spartan 3 FPGA} using System Verilog.
          \item Demonstrated recursion, looping and conditionals on the Computer by using a simple assembly language.
      \end{itemize}
    }

\scproject
    {Winter 2012-2013}
    {Conference Dates Web Crawler}
    {\textsc{Computing Laboratory}}
    {\href{https://github.com/srijanshetty/crawler} {\footnotesize{ [Repository]}} }
    {
           Developed a python based web crawler to crawl paper submission deadlines for a given paper using BeautifulSoup and
           urllib.

    }

%----------------------------------------------------------------------------------------
%  INDEPENDENT PROJECTS
%----------------------------------------------------------------------------------------
\section{Independent Projects}

\siproject {July 2014}
          {OARS}
          {
              Scraped the institute academic records to create a database of all courses offered by the institute to create an
              AngularJS based course search over the existing legacy course search.  \href{https://navya.github.io/oars}{[Website]}
           }

\siproject {March 2014}
          {Get Your Personal Homepage (GYPH)}
          {
               Eased the process of creating minimalistic websites for not-so-tech-savvy students by creating a WYSIWG website
               editor.  \href{http://gyph2.herokuapp.com/} {[Website]}
           }

\siproject {August 2013}
           {ShuffleRun}
           {
               Pitched ShuffleRun --- a music player which selects tracks on the basis of the user's running speed;
               recieved honourable mention at \textbf{Yahoo! HackU 2013}.  \href{https://github.com/srijanshetty/ShuffleRun} {[Repository]}
           }

\siproject {Summer 2012}
           {GNU/Linux Exploration, Programming Club, IIT Kanpur}
           {
              Explored and documented the various facets of \textbf{GNU/Linux} like the file system,
              process management, memory management, shell interface etc to serve as a reference for the
              Student Community.
              \href{https://docs.google.com/document/d/1ZHO9w36aoq3oaZBR4Um1AOmDfiTDAEgM6baQAu3icw4/edit?usp=sharing}{[Documentation]}
           }

\lproject {Summer 2012}
          {Voicing, Archiving and Networking the Information \textsc{(VANI)}, IIT Kanpur}
          {Professor Manindra Agarwal, Professor T. V. Prabhakar}
          {
              \begin{itemize}[leftmargin=0.5cm]
                  \item Conceptualized the idea of VANI - an extensible student community platform consisting
                      of a Student Wiki, Forums, Community Search, Calendar and Lost and Found.
                  \item Developed the Lost and Found module using \textbf{Drupal}
                      and conducted sessions on the extending VANI to create rich applications.
              \end{itemize}
          }

\projectlist {Projects on Github}
             {
                 \begin{itemize}[leftmargin=0.5cm]
                     \item \textbf{DS}: implementation of certain algorithms and data structures
                         \href{https://github.com/srijanshetty/DS} {https://github.com/srijanshetty/DS}
                     \item \textbf{Notes}: notes on computing and social sciences.
                         \href{https://github.com/srijanshetty/notes} {https://github.com/srijanshetty/notes}
                     \item \textbf{Dotfiles}: an opinionated work flow on Linux Systems.
                         \href{https://github.com/srijanshetty/dotfiles} {https://github.com/srijanshetty/dotfiles}
                     \item \textbf{Bash Scripts}: handy configurations and shortcuts for ZSH
                         \href{https://github.com/srijanshetty/custom} {https://github.com/srijanshetty/custom}
                 \end{itemize}
             }

\projectlist {Web Development}
             {
                 \begin{itemize}[leftmargin=0.5cm]
                     \item \textbf{Udghosh}: The annual sports-fest of IIT Kanpur.
                         \href{https://github.com/srijanshetty/udghosh-final}{https://github.com/srijanshetty/udghosh-final}
                     \item \textbf{ALI}: Antaragni Leadership Initiative.
                         \href{http://www.antargni.in/ali} {http://www.antargni.in/ali}
                 \end{itemize}
             }


%----------------------------------------------------------------------------------------
%   OTHER SKILLS
%----------------------------------------------------------------------------------------

\section {Social, Leadership and Artistic Skills}

\skill {\href{srijanshetty.quora.com} {[Blog]}}
       {Active blogger at \href{srijanshetty.quora.com} {srijanshetty.quora.com}}

\iproject {\href{https://github.com/navya} {[Github]} }
          {Hacker, Navya, FOSS Group IIT Kanpur}
          {
               \begin{itemize}[leftmargin=0.5cm]
                   \item Navya is the resident Free and Open Source Software Group of IIT Kanpur.
                   \item Promoted the FOSS initiative in campus by organizing student lectures and meet-ups.
                   \item Built student-centric applications like Course Search and Student Search, to circumvent
                       their legacy institute counterparts.
               \end{itemize}
          }

\iproject {2014}
          {Head, Major Events and Competitions, Antaragni '14}
          {
               \begin{itemize}[leftmargin=0.5cm]
                    \item Orchestrated Roadtrip --- the National Campaign of Antaragni '14, under which national prelims for Nukkad,
                        Dance, Quiz and Synchronicity were held in Pune, Mumbai, Chennai, Bangalore, Chandigarh, Bhopal, Lucknow,
                        Kolkata, and Kanpur.
                    \item Negotiated with leading academies of Dance, Musicals, and Fine Arts to provide mentorships, and other
                        non-monetary incentives to the outstanding cultural performers in Antaragni under the banner
                        of \textbf{Dream On} campaign.
                    \item Coordinated the conduction of fifty competitions and nine major events with a team comprising of over
                        thousand volunteers.
                    \item Initiated and organized Mr. \& Ms. Fresher's; and Choreo Nite to the arsenal of competitions
                        held in Antaragni.
                    \item Publicized Antaragni on social media platforms by contributing creative content.
               \end{itemize}
          }

\iproject {2013 -- 2014}
          {Hospitality Coordinator, Antaragni '13}
          {
               \begin{itemize}[leftmargin=0.5cm]
                   \item Overhauled the methodology of inviting colleges all around the country by calling college societies over college cultural unions.
                   \item Chalked out the logistics of the hospitality and accommodation of \textbf{1500 participants} with my fellow coordinators.
                   \item Helmed a team of 40 secretaries and 80 volunteers tasked with ensuring a flawless conduction of the festival.
                   \item Handled security during the four days of the Festival from 24$^{th}$ to 27$^{th}$ October.
               \end{itemize}
          }

\iproject {2013 -- 2014}
          {Member, Gymkhana Review Committee (GRC)}
          {
               \begin{itemize}[leftmargin=0.5cm]
                   \item The GRC was established to revamp the Students' Gymkhana of IIT Kanpur,
                       which is responsible for all the student activities of the institute.
                   \item Chaired the meetings on \textbf{Extended Orientation of Incoming Freshmen}.
                   \item Contributed actively as a member of academic, senate and activities sub-committees of the GRC.
               \end{itemize}
          }

\skill {2013 -- 2014}
       {Editor-in-chief, Vox Populi, the campus newsletter.}

\skill {2012 -- 2013}
       {Senator, Students' Senate, IITK Y11 batch.}

\skill {2010 -- 2011}
       {Basic training in \textbf{Indian Classical Music}. (Prayag Sangeet Samiti Allahabad).}

\skill {2005 -- 2006}
       {House Captain, Einstein House.}

%----------------------------------------------------------------------------------------
%   COMMUNITY SERVICE
%----------------------------------------------------------------------------------------

\section{Volunteer Work}

\iproject {2014 -- 2015}
          {Member, Institute Anti Ragging Committee}
          {
               \begin{itemize}[leftmargin=0.5cm]
                   \item Appointed as a member of the Institute Anti Ragging Committee, a 10 member committee responsible
                       for welfare of the newly admitted freshmen.
                   \item Worked in tandem with the office of the Dean of Student Affairs, to organize sessions
                       on easing the transition of freshmen to college life.
                   \item Ensured institute norms regarding humane treatment of freshmen are properly conveyed to all senior batches
                       and worked towards the enforcement of those norms.
               \end{itemize}
          }

\lproject {2013 -- 2014}
          {Community Service, English Proficiency Programme}
          {Professor Bhaskar Das Gupta}
          {
               \begin{itemize}[leftmargin=0.5cm]
                   \item A brainchild of \textbf{Professor Bhaskar Das Gupta}, English Proficiency Programme
                       tries to impart a functional knowledge of English to students.
                   \item Worked as a English tutor in the pilot programme to help under-privileged students from
                       and around IIT Kanpur campus; and helped in the conduction of second phase of the project
                       aimed at school teachers in nearby schools.
               \end{itemize}
          }

\lproject {2012 -- 2013}
          {Academic Mentor, Fundamentals of Computing}
          {Professor Sumit Ganguly}
          {
               \begin{itemize}[leftmargin=0.5cm]
                   \item Tutored four students in Fundamentals of Computing, a beginner level course on C.
                   \item Aided the instructor in understanding the difficulties faced by students by acting as a link
                       student between the students and instructor.
               \end{itemize}
          }

\iproject {2014 -- 2015}
          {Social Initiative, Antaragni '14}
          {
               \begin{itemize}[leftmargin=0.5cm]
                   \item Organized a blood donation camp for the residents of IIT Kanpur.
                   \item Mobilized information about road safety through infographics and questionnaires
                       via Antaragni's nation wide social campaign \textbf{Zara Sambhal Ke}.
                   \item Organized a Respect March on the 15$^{th}$ of August, to infuse a sense of patriotism
                       on the occassion of Indian Independence day.
               \end{itemize}
          }

\end{document}
